\documentclass[11pt]{amsart}
\usepackage[a4paper,margin=3cm]{geometry}
\usepackage{amsmath,amssymb,amsthm}
\usepackage{enumitem}
\usepackage{microtype}
\usepackage[colorlinks=true,linkcolor=blue,citecolor=blue,urlcolor=blue]{hyperref}

\newtheorem{theorem}{Theorem}
\newtheorem{lemma}{Lemma}[section]
\newtheorem{proposition}[lemma]{Proposition}
\theoremstyle{remark}
\newtheorem{remark}[lemma]{Remark}
\DeclareMathOperator{\Aut}{Aut}
\DeclareMathOperator{\Sym}{Sym}
\DeclareMathOperator{\im}{im}
\DeclareMathOperator{\rk}{rk}
\newcommand{\Q}{\mathbb Q}
\newcommand{\CC}{\mathbb C}
\newcommand{\PP}{\mathbb P}
\newcommand{\LL}{\mathcal L}
\numberwithin{equation}{section}
\allowdisplaybreaks[1]
\setlength{\parindent}{15pt}
\setlength{\parskip}{0pt plus 1pt}

\title[The tangent Gauss map]{A counterexample to generic injectivity of the tangent Gauss map}
\author{Xiaoxiang Zhou}
\address{Institut f\"ur Mathematik, Humboldt-Universit\"at zu Berlin,
12489 Berlin, Germany}
\email{xiaoxiang.zhou@hu-berlin.de}
\date{November 3, 2026}

\begin{document}
\maketitle

We work over $\CC$. For an irreducible curve $C$ in an abelian variety $A$,
its \emph{tangent Gauss map} is
\[
\phi_C:C^{\mathrm{sm}}\longrightarrow\PP(T_0A),\qquad
p\longmapsto d(t_{-p})_p(T_pC),
\]
where $t_a$ denotes translation by $a$. We call $C$ \emph{nondegenerate}
if it is not contained in a translate of a proper abelian subvariety.
Every automorphism of $A$ has the form $t_a\circ u$, with
$u\in\Aut_{\mathrm{gp}}(A)$; its differential is called scalar if
$du_0=\lambda\,\mathrm{id}_{T_0A}$ for some $\lambda\in\CC^\times$.

\begin{theorem}\label{thm:main}
There exist a simple abelian threefold $A$ and a nondegenerate irreducible
curve $C\subset A$ whose tangent Gauss map has generic degree $3$, but no
nonidentity automorphism of $A$ with scalar differential preserves $C$.
\end{theorem}

\section{Strategy}

\subsection{A pencil whose square gives the Gauss map}
Let $X$ be a smooth nonhyperelliptic curve of genus $4$, and let $\LL$ be
a line bundle of degree $3$ with $h^0(X,\LL)=2$. It is base-point-free:
otherwise, removing a base point would give a degree-two pencil, making
$X$ hyperelliptic. Thus $\LL$ defines a degree-three map
\[
\psi_{\LL}:X\longrightarrow\PP^1.
\]
The products of its sections define a map to $\PP^2$ that factors through
$\psi_{\LL}$. In a basis $H^0(X,\LL)=\langle s_0,s_1\rangle$, this map is
\begin{equation}\label{eq:veronese}
\gamma=[s_0^2:s_0s_1:s_1^2]=\nu_2\circ\psi_{\LL},
\end{equation}
where $\nu_2:\PP^1\hookrightarrow\PP^2$ is the quadratic Veronese
embedding. In particular, $\gamma$ has degree $3$ onto its image.

Our aim is to realize this map as a tangent Gauss map. For this, we seek
the following data:
\begin{enumerate}[label=\arabic*.,leftmargin=*]
\item a theta characteristic $\LL$, that is,
$\LL^{\otimes2}\cong\omega_X$, with $h^0(X,\LL)=2$;
\item a quotient $q:J(X)\twoheadrightarrow A$ onto a simple abelian
threefold such that
\begin{equation}\label{eq:quotient}
q^*H^0(A,\Omega_A^1)=W(X,\LL),\qquad
W(X,\LL):=\im\!\left(
\Sym^2H^0(X,\LL)\longrightarrow H^0(X,\omega_X)\right).
\end{equation}
\end{enumerate}
Here we identify invariant differentials on $J(X)$ with canonical forms
on $X$. We say that $W(X,\LL)$ \emph{comes from an abelian quotient}
when the equality in \eqref{eq:quotient} holds, without requiring the
quotient to be simple.

For any two independent sections of a line bundle, their three quadratic
products are linearly independent: a relation would give a constant
quadratic equation for their nonconstant ratio. Thus $W(X,\LL)$ has
dimension $3$ whenever $h^0(X,\LL)=2$.

If \eqref{eq:quotient} holds, composing an Abel--Jacobi map with $q$ gives
$f:X\to A$. The pullbacks of invariant forms are $s_0^2,s_0s_1,s_1^2$,
so its tangent map is precisely \eqref{eq:veronese}. In Section~\ref{sec:completion}
we show that $X$ is the normalization of $C=f(X)$ and exclude the
automorphisms in Theorem~\ref{thm:main}.

\subsection{Starting from a hyperelliptic curve}
The two requirements above are easier to arrange on a hyperelliptic
curve. Recall that a finite map $\pi:(X,p)\to(E,0)$ to an elliptic curve
is \emph{tangential at $p$} if the elliptic subvariety
$\pi^*E\subset J(X)$ is tangent to the Abel--Jacobi image of $X$ at the
origin. We identify $E$ with its dual by its canonical principal
polarization.

\begin{lemma}\label{lem:tangency}
Let $X$ be a smooth hyperelliptic curve of genus at least $2$, and let
$p\in X$ be a Weierstrass point. The subspace
$H^0(X,\omega_X(-2p))$ comes from an abelian quotient if and only if
there is a tangential cover $\pi:(X,p)\to(E,0)$.
\end{lemma}

\begin{proof}
At a Weierstrass point,
\[
H^0(X,\omega_X(-p))=H^0(X,\omega_X(-2p)).
\]
Indeed, both spaces have dimension $g(X)-1$, by Riemann--Roch and
$h^0(X,\mathcal O_X(2p))=2$. Under the identification
$T_0^*J(X)=H^0(X,\omega_X)$, this space is the annihilator of the tangent
line to the Abel--Jacobi image at $p$.

If $\pi$ is tangential, take the quotient
$q:J(X)\to J(X)/\pi^*E$. Its pulled-back cotangent space is exactly this
annihilator. Conversely, suppose a quotient has this pulled-back
cotangent space. Its connected kernel $K$ is an elliptic subvariety
whose tangent line agrees with that of $X$ at $p$. Dualizing the
inclusion $K\hookrightarrow J(X)$, using the principal polarization of
$J(X)$, and composing with the Abel--Jacobi map based at $p$, gives a
nonconstant map $X\to K^\vee$. Its dual elliptic image is $K$, so the
map is tangential at $p$.
\end{proof}

\begin{lemma}[The seed]\label{lem:seed}
There exist a smooth hyperelliptic curve $X_0$ of genus $4$, a
Weierstrass point $p_0\in X_0$, and a tangential cover
$\pi_0:(X_0,p_0)\to(E_0,0)$ such that
\[
A_0=J(X_0)/\pi_0^*E_0
\]
is a simple abelian threefold.
\end{lemma}

The construction and the arithmetic check of simplicity are given in
Appendix~\ref{app:seed}. Fix such a seed and put
\[
\LL_0=\mathcal O_{X_0}(3p_0),\qquad
q_0:J(X_0)\twoheadrightarrow A_0.
\]
The hyperelliptic gap sequence gives
$\omega_{X_0}\cong\mathcal O_{X_0}(6p_0)$ and $h^0(X_0,\LL_0)=2$.
If $x$ is a hyperelliptic coordinate with a double pole at $p_0$, the
sections of $\LL_0$ are spanned by $1,x$, and their products by
$1,x,x^2$. Therefore
\begin{equation}\label{eq:seed-W}
\begin{split}
W_0:=W(X_0,\LL_0)
&=H^0(X_0,\omega_{X_0}(-2p_0))\\
&=q_0^*H^0(A_0,\Omega_{A_0}^1).
\end{split}
\end{equation}
Thus the seed has the required theta characteristic and quotient.
We now deform it while preserving these properties and leaving the
hyperelliptic locus.

\section{A dimension count for deformations}\label{sec:deformation}

Let $B$ be a sufficiently small, simply connected, marked local
deformation space of the spin curve $(X_0,\LL_0)$. Write
$(X_t,\LL_t)$ for the corresponding family. Since theta characteristics
deform discretely with the curve,
\[
\dim B=h^1(X_0,\mathcal T_{X_0})=3g-3=9.
\]
We impose the two requirements in Section~1 and then estimate the
dimension of the hyperelliptic part.

\subsection{The theta-null condition}
Let
\[
S=\{t\in B:h^0(X_t,\LL_t)\geq2\}.
\]
The spin structure is even. Locally, $S$ is the zero locus of its theta
constant, so every component through the seed has dimension at least
$8$. After shrinking $B$, upper semicontinuity and parity imply
$h^0(X_t,\LL_t)=2$ on $S$. Cohomology and base change, together with the
injectivity of multiplication noted above, then show that
\[
W_t=W(X_t,\LL_t)
\]
forms a holomorphic rank-three subbundle of the Hodge bundle over $S$.
We use the reduced analytic structures on these loci.

\subsection{Preserving the abelian quotient}
Define the rational six-dimensional subspace
\[
V_0=q_0^*H^1(A_0,\Q)\subset H^1(X_0,\Q).
\]
Parallel transport under the Gauss--Manin connection gives
$V_t\subset H^1(X_t,\Q)$. The marking identifies these spaces with a
fixed rational vector space. By \eqref{eq:seed-W},
\[
V_{0,\CC}=W_0\oplus\overline{W_0}.
\]

For $t\in S$, we have the equivalences
\begin{equation}\label{eq:hodge}
\begin{aligned}
&\left\{\begin{array}{l}
V_t\text{ is a Hodge substructure},\\
V_{t,\CC}\cap H^0(X_t,\omega_{X_t})=W_t
\end{array}\right.\\
&\hspace{1cm}\Longleftrightarrow\quad
V_{t,\CC}=W_t\oplus\overline{W_t}\\
&\hspace{1cm}\Longleftrightarrow\quad W_t\subset V_{t,\CC}\\
&\hspace{1cm}\Longleftrightarrow\quad
\rk\!\left(W_t\longrightarrow H^1(X_t,\CC)/V_{t,\CC}\right)=0.
\end{aligned}
\end{equation}
For the middle equivalence, rationality makes $V_{t,\CC}$ stable under
complex conjugation. The two spaces $W_t$ and $\overline{W_t}$ have
dimension $3$ and intersect trivially by Hodge type, so their direct
sum has the full dimension $6$ of $V_{t,\CC}$.

Let
\[
Z=\{t\in S:W_t\subset V_{t,\CC}\}.
\]
It is the zero locus of a holomorphic bundle map of ranks $3$ and $2$,
hence is locally defined in $S$ by at most six equations. Every
irreducible component of $Z$ through the seed therefore satisfies
\begin{equation}\label{eq:lower-bound}
\dim Z\geq 9-1-6=2.
\end{equation}
From now on, fix one such irreducible component and denote it again by
$Z$.

On $Z$, the space $V_t$ is a polarizable rational Hodge substructure.
The correspondence between such weight-one Hodge structures and
abelian varieties up to isogeny gives a quotient
\[
q_t:J(X_t)\twoheadrightarrow A_t,\qquad
q_t^*H^0(A_t,\Omega_{A_t}^1)=W_t.
\]
We may take $q_t$ to have connected kernel by choosing the corresponding
quotient of $J(X_t)$. Thus $W_t$ continues to come from an abelian
quotient.

\subsection{Leaving the hyperelliptic locus}
Write $\mathcal H_4\subset B$ for the hyperelliptic locus. If
$t\in S\cap\mathcal H_4$, let $H_t$ be the hyperelliptic pencil on
$X_t$. A degree-three line bundle with two sections on such a curve
has the form $H_t(p_t)$. Indeed, a base-point-free degree-three pencil,
together with the degree-two pencil, would give a birational map to
a curve of bidegree $(2,3)$ in $\PP^1\times\PP^1$, whose arithmetic
genus is $2$, contradicting $g(X_t)=4$. Removing the base point leaves
the hyperelliptic pencil.

Now use the theta-characteristic condition. Since
$\omega_{X_t}\cong H_t^{\otimes3}$, the equality
$\LL_t^{\otimes2}\cong\omega_{X_t}$ gives
$\mathcal O_{X_t}(2p_t)\cong H_t$. Thus $p_t$ is a Weierstrass point and
\begin{equation}\label{eq:hyperelliptic-W}
\LL_t\cong\mathcal O_{X_t}(3p_t),\qquad
W_t=H^0(X_t,\omega_{X_t}(-2p_t)).
\end{equation}
There are only finitely many choices of the marked Weierstrass point;
locally we can work on each marked branch separately.

\begin{lemma}[Finiteness of tangential covers]\label{lem:finiteness}
Fix an elliptic curve $E$ and a positive integer $n$. Up to isomorphism,
there are only finitely many smooth hyperelliptic tangential covers
$\pi:(X,p)\to(E,0)$ of degree $n$.
\end{lemma}

\begin{proof}
By Treibich's correspondence, these covers are represented by rational
irreducible curves in linear systems indexed by
$\mu\in\mathbb Z_{\geq0}^4$ and $d\geq0$, with
\[
2n+1=\sum_{i=0}^3\mu_i^2+4d.
\]
For fixed $n$, only finitely many such indices occur. Each corresponding
Severi set is finite by \cite[Proposition~2.5, 3.]{treibich2025}.
The correspondence is described in \cite[Section~2, Remark~1, 3.--4.]{treibich2025}.
\end{proof}

To apply this lemma, we must keep the degree fixed. On $Z$, let
$K_t=(\ker q_t)^0$, an elliptic subvariety of $J(X_t)$. Its integral
homology lattice is the annihilator of $V_t$ in $H_1(X_t,\mathbb Z)$,
so it is fixed under the marking. Locally these kernels form a family
of elliptic curves. Dualizing their inclusions using the principal
polarizations of the Jacobians gives maps
\[
\pi_t:(X_t,p_t)\longrightarrow(E_t,0),\qquad E_t=K_t^\vee.
\]
The induced integral cohomology maps are fixed under local markings.
Their cup-product pairings therefore show that $\deg\pi_t=n$ is
locally constant. By Lemma~\ref{lem:tangency} and
\eqref{eq:hyperelliptic-W}, these maps are tangential on the
hyperelliptic branch.

Lemma~\ref{lem:finiteness} gives finite fibers over each elliptic
isomorphism class. Working in a local deformation chart introduces
only finite choices of spin structure and marked Weierstrass point.
Since elliptic curves have a one-dimensional moduli space, it follows
that
\begin{equation}\label{eq:upper-bound}
\dim(Z\cap\mathcal H_4)\leq1.
\end{equation}
Together, \eqref{eq:lower-bound} and \eqref{eq:upper-bound} show that
$Z$ contains nonhyperelliptic points.

\subsection{Choosing a simple quotient}
It remains to ensure that the quotient at a chosen nonhyperelliptic
point is simple. The cup-product polarization restricts to a fixed
nondegenerate rational form on $V_0$. For each proper nonzero rational
subspace $R\subset V_0$ on which this form is nondegenerate, let $e_R$
be the orthogonal projector onto $R$. The locus
\[
Z_R=\{t\in Z:e_R(W_t)\subset W_t\}
\]
is closed analytic. It is exactly the locus where the transported
subspace $R$ is a Hodge substructure. It is proper, since the seed
quotient $A_0$ is simple. Any Hodge substructure is nondegenerate for
the polarization, so these are all the subspaces that need to be
excluded.

There are only countably many rational subspaces $R$. By the Baire
category theorem, an irreducible complex analytic space cannot be
covered by countably many proper closed analytic subsets. We may
therefore choose
\[
t\in Z\setminus\left(\mathcal H_4\cup\bigcup_R Z_R\right).
\]
For this point, $X=X_t$ is nonhyperelliptic and $A=A_t$ is simple.
With $\LL=\LL_t$ and $q=q_t$, we have obtained both requirements of
Section~1.

\section{Completing the proof}\label{sec:completion}

\subsection{Passing to the image curve}
Choose a base point $b\in X$ and set
\[
f=q\circ j_b:X\longrightarrow A,\qquad C=f(X),
\]
where $j_b$ is the Abel--Jacobi map. The forms $s_0^2,s_0s_1,s_1^2$
have no common zero, so $f$ is immersive. Moreover, $C$ contains the
origin and generates $A$, because $j_b(X)$ generates $J(X)$ and $q$ is
surjective. In particular, $C$ is nondegenerate.

Let $Y$ be the normalization of $C$, and factor $f$ through a finite map
$r:X\to Y$ of degree $e$. The tangent map $\gamma$ in
\eqref{eq:veronese} factors through $r$, since it is the pullback of
the Gauss map on the smooth locus of $C$. Thus $e$ divides $3$. If
$e>1$, then $e=3$ and the induced map from $Y$ to the conic
$\nu_2(\PP^1)$ has degree $1$. This would make $Y\cong\PP^1$, which
is impossible: an abelian variety contains no nonconstant rational
curve. Hence $e=1$, and $X$ is the normalization of $C$. Consequently,
$\phi_C$ has generic degree $3$.

\subsection{Excluding scalar symmetries}
\begin{lemma}\label{lem:scalar}
Let $A$ be a simple complex abelian variety of dimension greater than
$1$. If $u\in\Aut_{\mathrm{gp}}(A)$ has scalar differential, then
$u=[1]$ or $u=[-1]$.
\end{lemma}

\begin{proof}
Write $du_0=\lambda\,\mathrm{id}$. On $H_1(A,\CC)$, the operator $u_*$
is diagonalizable with eigenvalues among $\lambda,\overline\lambda$.
Since it is rational, the minimal polynomial of $\lambda$ over $\Q$
has degree at most $2$.

If $\lambda\notin\mathbb R$, then $K=\Q(\lambda)$ is imaginary
quadratic. Write $\lambda=a+ib$, with $b\neq0$. On $H_1(A,\mathbb R)$,
the complex structure is
\[
J=\frac{u_*-a\,\mathrm{id}}{b}.
\]
Thus every one-dimensional $K$-subspace of $H_1(A,\Q)$ has a
$J$-invariant realification. Its intersection with the integral lattice
defines an elliptic subvariety of $A$, contradicting simplicity.

Hence $\lambda\in\mathbb R$. Taking the trace gives $\lambda\in\Q$.
Both $u_*$ and its inverse preserve the integral lattice, so $\lambda$
is a rational algebraic unit and therefore equals $1$ or $-1$. A
homomorphism of complex abelian varieties is determined by its
differential, proving the assertion.
\end{proof}

\begin{proof}[Conclusion of the proof of Theorem~\ref{thm:main}]
Suppose that $\tau=t_a\circ u$ preserves $C$ and has scalar
differential. It lifts to an automorphism $\sigma$ of the normalization
$X$. If $\sigma$ is the identity, then $\tau$ fixes $C$ pointwise.
Since $0\in C$ and $C$ generates $A$, this forces $a=0$ and $u=[1]$.
We may therefore assume that $\sigma$ is nontrivial.

A scalar differential preserves every tangent line, so
\[
\gamma\circ\sigma=\gamma.
\]
Lemma~\ref{lem:scalar} gives $u=[1]$ or $u=[-1]$. If $u=[-1]$, then
$\tau$ is an involution, and its nontrivial lift $\sigma$ has order
$2$. The degree-three map $\gamma$ would factor through
$X/\langle\sigma\rangle$, forcing $2\mid3$, a contradiction.

If $u=[1]$, then $\tau=t_a$ with $a\neq0$. Since $g(X)=4$, the
automorphism $\sigma$ has finite order $m$. Its cyclic group acts
freely on a generic fiber of $\gamma$, so $m\mid3$ and hence $m=3$.
The induced map $X/\langle\sigma\rangle\to\nu_2(\PP^1)$ has degree
$1$, giving $X/\langle\sigma\rangle\cong\PP^1$. On the other hand,
$\sigma$ has no fixed point, because
$f(\sigma(x))=f(x)+a$ and a nonzero translation has none.
Riemann--Hurwitz for this unramified degree-three quotient would give
\[
6=2g(X)-2=3\bigl(2g(\PP^1)-2\bigr)=-6,
\]
again a contradiction. Thus $\tau$ is the identity.
\end{proof}

\begin{remark}
The construction gives an immersive normalization $X\to C$. It does
not require the image curve $C$ to be smooth or the map $X\to A$ to be
an embedding.
\end{remark}

\appendix
\section{A seed with a simple complementary quotient}\label{app:seed}

We construct the seed of Lemma~\ref{lem:seed} from a genus-four
Lam\'e curve. Reduction modulo $11$ will verify smoothness and
simplicity.

\subsection{The tangential cover}
Set
\[
E_0:y^2=4x^3-4x-1,
\]
and let $e_1,e_2,e_3$ be the roots of $4e^3-4e-1$. Define
\[
L(B)=(B^3-208B+560)
\prod_{i=1}^3(B^2+10e_iB-35e_i^2-28).
\]
This is the genus-four Lam\'e spectral polynomial for $g_2=4$, $g_3=1$;
see \cite[Table~3 and Theorem~L]{maier}. The associated formulas give
a degree-ten cover from the smooth projective model of
$\nu^2=L(B)$ to $E_0$, sending the point at infinity to the origin.
The cover is tangential at infinity: the Lam\'e operator has an
odd-order commuting operator with elliptic coefficients, and its
spectral cover is tangential by \cite[Proposition~3.1 and the paragraph
after equation~(17)]{eilbeck}.

The polynomial and cover are defined over $\Q$. The change of variable
$Y=4\nu$ gives the integral model $Y^2=16L(B)$. The Weierstrass
discriminant of $E_0$ is $37$, and reduction of $L$ modulo $11$ gives
\[
\overline L=B^9+B^4+6B^3+9B^2+7B+5
\]
with factorization
\[
\overline L=(B+2)(B^2+9B+5)
(B^3+5B^2+6B+1)(B^3+6B^2+7B+6).
\]
These are distinct irreducible factors over $\mathbb F_{11}$. Thus
the curve $X_0$ is smooth of genus $4$ and has good reduction at $11$.
Its unique point at infinity, denoted $p_0$, is a Weierstrass point.
This supplies the required tangential cover $\pi_0$.

\subsection{Checking simplicity}
Let $N_r=\#X_0(\mathbb F_{11^r})$. Direct point counting gives
\[
(N_1,N_2,N_3,N_4)=(10,120,1312,15312).
\]
These numbers can be reproduced from
\[
N_r=11^r+1+\sum_{b\in\mathbb F_{11^r}}\chi_r(\overline L(b)),
\]
where $\chi_r$ is the quadratic character extended by $\chi_r(0)=0$.
For $r=2,3,4$, one may use the respective irreducible defining
polynomials
\[
z^2+1,\qquad z^3+4z^2+1,\qquad z^4+4z^3+1.
\]
The four character sums are $-2,-2,-20,670$. Newton's identities and
the functional equation give the Frobenius polynomial
\[
\begin{split}
P_{X_0,11}(T)={}&T^8-2T^7+T^6-6T^5+180T^4\\
&-66T^3+121T^2-2662T+14641.
\end{split}
\]
Also, $\#E_0(\mathbb F_{11})=7$, so
$P_{E_0,11}(T)=T^2-5T+11$.

The maps on Jacobians extend across good reduction. Since
$\pi_{0,*}\pi_0^*=[10]$ and $11\nmid10$, the elliptic factor remains
an elliptic factor in the special fiber. The quotient
$A_0=J(X_0)/\pi_0^*E_0$ has good reduction, with Frobenius polynomial
\[
\begin{split}
h(T)&=P_{X_0,11}(T)/P_{E_0,11}(T)\\
&=T^6+3T^5+5T^4-14T^3+55T^2+363T+1331.
\end{split}
\]
Modulo $3$, this becomes
$\overline h=T^6-T^4+T^3+T^2-1$. The finite-field irreducibility test
gives
\[
\gcd(\overline h,T^9-T)=\gcd(\overline h,T^{27}-T)=1,
\qquad T^{729}\equiv T\pmod{\overline h}.
\]
Hence $h$ is irreducible over $\Q$, and the special fiber of $A_0$ is
simple. It is ordinary because $11\nmid14$. By
\cite[Proposition~3.4]{arias}, an ordinary simple abelian threefold
over a finite field of odd cardinality is absolutely simple unless
its Frobenius polynomial has the form $T^6+cT^3+q^3$. Our polynomial
does not have this form, so the special fiber is absolutely simple.

Absolute simplicity then holds in characteristic zero: a proper
positive-dimensional geometric abelian subvariety would, after a
finite base extension, specialize to such a subvariety in the good
special fiber. Thus $A_0$ is simple over $\CC$, proving
Lemma~\ref{lem:seed}.

The accompanying script \texttt{verify\_tangent\_seed.py} reproduces
the polynomial reduction, point counts, Frobenius factorization, and
irreducibility checks using only the Python standard library.

\begin{thebibliography}{9}
\bibitem{arias}
S.~Arias-de-Reyna, C.~Armana, V.~Karemaker, M.~Rebolledo,
L.~Thomas, and N.~Vila,
\emph{Large Galois images for Jacobian varieties of genus 3 curves},
2016, \href{https://arxiv.org/abs/1507.05913}{arXiv:1507.05913}.

\bibitem{eilbeck}
J.~C.~Eilbeck, V.~Z.~Enolskii, and E.~Previato,
\emph{Varieties of elliptic solitons},
J.\ Phys.\ A: Math.\ Gen.\ \textbf{34} (2001), 2215--2227,
\href{https://doi.org/10.1088/0305-4470/34/11/314}{doi:10.1088/0305-4470/34/11/314}.

\bibitem{maier}
R.~S.~Maier,
\emph{Lam\'e polynomials, hyperelliptic reductions and Lam\'e band structure},
\href{https://arxiv.org/abs/math-ph/0309005}{arXiv:math-ph/0309005},
version~4, 2004.

\bibitem{treibich2025}
A.~Treibich,
\emph{Hyperelliptic tangential covers and even elliptic finite-gap
potentials, back and forth},
\href{https://arxiv.org/abs/2501.16483v1}{arXiv:2501.16483v1}, 2025.
\end{thebibliography}
\end{document}
